%%%%%%%%%%%%%%%%%%%%%%%%%%%%%%%%%%%%%%%%%%%%%%%%%%%%%%%%%%%%
% 7.12 STOCHASTIC PROCESSES
% The Composition of Advanced Mathematics
%%%%%%%%%%%%%%%%%%%%%%%%%%%%%%%%%%%%%%%%%%%%%%%%%%%%%%%%%%%%

\section{Stochastic Processes}
\label{sec:stochastic-processes}

\prerequisite{
Random variables, joint distributions, conditional probability,
expectation, covariance, measure-theoretic probability, filtrations,
conditional expectation, laws of large numbers, and central limit
theorems.
}

\learningobjectives{
By the end of this section, the reader should be able to:
\begin{itemize}
    \item define a stochastic process;
    \item distinguish discrete-time and continuous-time processes;
    \item distinguish discrete-state and continuous-state processes;
    \item interpret sample paths;
    \item define finite-dimensional distributions;
    \item understand the role of filtrations;
    \item define adapted processes;
    \item define stationary and weakly stationary processes;
    \item define independent increments;
    \item define stationary increments;
    \item work with Bernoulli processes;
    \item work with simple random walks;
    \item compute random-walk means and variances;
    \item understand first-passage and hitting times;
    \item define Markov chains;
    \item work with transition matrices;
    \item compute multi-step transition probabilities;
    \item understand communicating classes and recurrence;
    \item define stationary distributions;
    \item understand detailed balance;
    \item define Poisson processes;
    \item derive Poisson increment distributions;
    \item understand exponential interarrival times;
    \item define counting processes;
    \item define renewal processes;
    \item understand branching processes;
    \item define martingales;
    \item identify basic martingale examples;
    \item define stopping times;
    \item understand optional-stopping ideas;
    \item define covariance functions;
    \item recognize Gaussian processes;
    \item recognize Brownian motion;
    \item connect stochastic processes with queues, reliability,
          finance, population models, signal processing, and dynamical
          systems.
\end{itemize}
}

%%%%%%%%%%%%%%%%%%%%%%%%%%%%%%%%%%%%%%%%%%%%%%%%%%%%%%%%%%%%
\subsection{What Is a Stochastic Process?}
\label{subsec:what-is-stochastic-process}
%%%%%%%%%%%%%%%%%%%%%%%%%%%%%%%%%%%%%%%%%%%%%%%%%%%%%%%%%%%%

A stochastic process is a collection of random variables indexed by a
parameter such as time or space.

\begin{definitionbox}{Stochastic Process}
A stochastic process is a family

\[
\boxed{
\{X_t:t\in T\},
}
\]

where each

\[
X_t
\]

is a random variable defined on a common probability space.

The set \(T\) is called the index set.
\end{definitionbox}

Often \(t\) represents time.

%%%%%%%%%%%%%%%%%%%%%%%%%%%%%%%%%%%%%%%%%%%%%%%%%%%%%%%%%%%%
\subsection{Discrete-Time and Continuous-Time Processes}
\label{subsec:discrete-continuous-time-processes}
%%%%%%%%%%%%%%%%%%%%%%%%%%%%%%%%%%%%%%%%%%%%%%%%%%%%%%%%%%%%

If

\[
T
=
\{0,1,2,\ldots\},
\]

the process is called discrete-time.

We may write

\[
\boxed{
X_0,X_1,X_2,\ldots
}
\]

or

\[
\boxed{
(X_n)_{n\geq0}.
}
\]

If

\[
T=[0,\infty)
\]

or

\[
T=\mathbb{R},
\]

the process is continuous-time.

We then often write

\[
\boxed{
(X_t)_{t\geq0}.
}
\]

%%%%%%%%%%%%%%%%%%%%%%%%%%%%%%%%%%%%%%%%%%%%%%%%%%%%%%%%%%%%
\subsection{State Space}
\label{subsec:stochastic-state-space}
%%%%%%%%%%%%%%%%%%%%%%%%%%%%%%%%%%%%%%%%%%%%%%%%%%%%%%%%%%%%

The set of values a stochastic process can take is called its state
space.

Examples include:

\[
\boxed{
\{0,1\}
}
\]

for a binary process,

\[
\boxed{
\{0,1,2,\ldots\}
}
\]

for a counting process,

and

\[
\boxed{
\mathbb{R}
}
\]

for a continuous-valued process.

Thus processes can be classified by both time and state space.

%%%%%%%%%%%%%%%%%%%%%%%%%%%%%%%%%%%%%%%%%%%%%%%%%%%%%%%%%%%%
\subsection{Four Basic Process Types}
\label{subsec:four-process-types}
%%%%%%%%%%%%%%%%%%%%%%%%%%%%%%%%%%%%%%%%%%%%%%%%%%%%%%%%%%%%

A useful classification is

\[
\begin{array}{c|c|c}
&
\text{Discrete State}
&
\text{Continuous State}
\\
\hline
\text{Discrete Time}
&
\text{Markov chain}
&
\text{random walk / time series}
\\[1mm]
\text{Continuous Time}
&
\text{Poisson process}
&
\text{Brownian motion}
\end{array}
\]

These examples are not exhaustive, but they illustrate the main
categories.

%%%%%%%%%%%%%%%%%%%%%%%%%%%%%%%%%%%%%%%%%%%%%%%%%%%%%%%%%%%%
\subsection{Sample Paths}
\label{subsec:sample-paths}
%%%%%%%%%%%%%%%%%%%%%%%%%%%%%%%%%%%%%%%%%%%%%%%%%%%%%%%%%%%%

Fix an outcome

\[
\omega\in\Omega.
\]

Then

\[
t\mapsto X_t(\omega)
\]

is an ordinary function of \(t\).

This function is called a sample path or realization of the process.

Thus a stochastic process can be viewed in two ways:

\[
\boxed{
t\mapsto X_t
}
\]

as a family of random variables,

or

\[
\boxed{
t\mapsto X_t(\omega)
}
\]

as a random path.

%%%%%%%%%%%%%%%%%%%%%%%%%%%%%%%%%%%%%%%%%%%%%%%%%%%%%%%%%%%%
\subsection{Finite-Dimensional Distributions}
\label{subsec:finite-dimensional-distributions}
%%%%%%%%%%%%%%%%%%%%%%%%%%%%%%%%%%%%%%%%%%%%%%%%%%%%%%%%%%%%

A stochastic process is described probabilistically through joint laws
such as

\[
(X_{t_1},X_{t_2},\ldots,X_{t_n}).
\]

These are called finite-dimensional distributions.

For times

\[
t_1<\cdots<t_n,
\]

we study probabilities of the form

\[
\boxed{
P
(
X_{t_1}\in A_1,
\ldots,
X_{t_n}\in A_n
).
}
\]

These joint distributions encode dependence across time.

%%%%%%%%%%%%%%%%%%%%%%%%%%%%%%%%%%%%%%%%%%%%%%%%%%%%%%%%%%%%
\subsection{Mean Function}
\label{subsec:process-mean-function}
%%%%%%%%%%%%%%%%%%%%%%%%%%%%%%%%%%%%%%%%%%%%%%%%%%%%%%%%%%%%

If expectations exist, define the mean function

\[
\boxed{
m_X(t)
=
\mathbb{E}[X_t].
}
\]

Unlike the mean of a single random variable, this may vary with time.

For example,

\[
m_X(t)=ct
\]

would describe linearly increasing expected behavior.

%%%%%%%%%%%%%%%%%%%%%%%%%%%%%%%%%%%%%%%%%%%%%%%%%%%%%%%%%%%%
\subsection{Covariance Function}
\label{subsec:process-covariance-function}
%%%%%%%%%%%%%%%%%%%%%%%%%%%%%%%%%%%%%%%%%%%%%%%%%%%%%%%%%%%%

If second moments exist, define

\[
\boxed{
C_X(s,t)
=
\operatorname{Cov}(X_s,X_t).
}
\]

Equivalently,

\[
\boxed{
C_X(s,t)
=
\mathbb{E}
[
(X_s-m_X(s))
(X_t-m_X(t))
].
}
\]

This measures dependence between process values at different times.

%%%%%%%%%%%%%%%%%%%%%%%%%%%%%%%%%%%%%%%%%%%%%%%%%%%%%%%%%%%%
\subsection{Autocovariance}
\label{subsec:autocovariance}
%%%%%%%%%%%%%%%%%%%%%%%%%%%%%%%%%%%%%%%%%%%%%%%%%%%%%%%%%%%%

When the process is stationary in an appropriate sense, covariance may
depend only on the lag

\[
h=t-s.
\]

Then one writes

\[
\boxed{
\gamma(h)
=
\operatorname{Cov}(X_t,X_{t+h}).
}
\]

The function

\[
\gamma
\]

is called the autocovariance function.

%%%%%%%%%%%%%%%%%%%%%%%%%%%%%%%%%%%%%%%%%%%%%%%%%%%%%%%%%%%%
\subsection{Autocorrelation}
\label{subsec:autocorrelation}
%%%%%%%%%%%%%%%%%%%%%%%%%%%%%%%%%%%%%%%%%%%%%%%%%%%%%%%%%%%%

If

\[
\gamma(0)>0,
\]

define the autocorrelation function

\[
\boxed{
\rho(h)
=
\frac{
\gamma(h)
}{
\gamma(0)
}.
}
\]

Thus,

\[
\boxed{
\rho(0)=1.
}
\]

Autocorrelation describes linear dependence across time lags.

%%%%%%%%%%%%%%%%%%%%%%%%%%%%%%%%%%%%%%%%%%%%%%%%%%%%%%%%%%%%
\subsection{Strict Stationarity}
\label{subsec:strict-stationarity}
%%%%%%%%%%%%%%%%%%%%%%%%%%%%%%%%%%%%%%%%%%%%%%%%%%%%%%%%%%%%

\begin{definitionbox}{Strictly Stationary Process}
A process \((X_t)\) is strictly stationary if for every \(n\), every
choice of times

\[
t_1,\ldots,t_n,
\]

and every shift \(h\),

\[
\boxed{
(X_{t_1},\ldots,X_{t_n})
\overset{d}{=}
(X_{t_1+h},\ldots,X_{t_n+h}).
}
\]
\end{definitionbox}

Thus all finite-dimensional distributions are invariant under time
translation.

%%%%%%%%%%%%%%%%%%%%%%%%%%%%%%%%%%%%%%%%%%%%%%%%%%%%%%%%%%%%
\subsection{Weak Stationarity}
\label{subsec:weak-stationarity}
%%%%%%%%%%%%%%%%%%%%%%%%%%%%%%%%%%%%%%%%%%%%%%%%%%%%%%%%%%%%

A process is weakly stationary, or second-order stationary, if:

\[
\boxed{
\mathbb{E}[X_t]=\mu
}
\]

does not depend on \(t\),

and

\[
\boxed{
\operatorname{Cov}(X_s,X_t)
=
\gamma(t-s)
}
\]

depends only on the lag.

Weak stationarity concerns only first and second moments.

%%%%%%%%%%%%%%%%%%%%%%%%%%%%%%%%%%%%%%%%%%%%%%%%%%%%%%%%%%%%
\subsection{Strict Versus Weak Stationarity}
\label{subsec:strict-versus-weak-stationarity}
%%%%%%%%%%%%%%%%%%%%%%%%%%%%%%%%%%%%%%%%%%%%%%%%%%%%%%%%%%%%

Strict stationarity concerns the full joint distribution.

Weak stationarity concerns only:

\[
\boxed{
\text{mean}
+
\text{covariance}.
}
\]

Strict stationarity does not automatically imply weak stationarity
unless the required second moments exist.

For Gaussian processes, weak stationarity is especially powerful
because the mean and covariance determine the full finite-dimensional
distributions.

%%%%%%%%%%%%%%%%%%%%%%%%%%%%%%%%%%%%%%%%%%%%%%%%%%%%%%%%%%%%
\subsection{Independent Increments}
\label{subsec:independent-increments}
%%%%%%%%%%%%%%%%%%%%%%%%%%%%%%%%%%%%%%%%%%%%%%%%%%%%%%%%%%%%

\begin{definitionbox}{Independent Increments}
A process \((X_t)\) has independent increments if increments over
disjoint time intervals are independent.
\end{definitionbox}

For

\[
0\leq t_0<t_1<\cdots<t_n,
\]

the increments

\[
X_{t_1}-X_{t_0},
\]

\[
X_{t_2}-X_{t_1},
\]

\[
\ldots,
\]

\[
X_{t_n}-X_{t_{n-1}}
\]

must be independent.

%%%%%%%%%%%%%%%%%%%%%%%%%%%%%%%%%%%%%%%%%%%%%%%%%%%%%%%%%%%%
\subsection{Stationary Increments}
\label{subsec:stationary-increments}
%%%%%%%%%%%%%%%%%%%%%%%%%%%%%%%%%%%%%%%%%%%%%%%%%%%%%%%%%%%%

A process has stationary increments if the distribution of

\[
X_{t+h}-X_t
\]

depends only on \(h\), not on \(t\).

Thus,

\[
\boxed{
X_{t+h}-X_t
\overset{d}{=}
X_h-X_0
}
\]

under suitable indexing.

Poisson processes and Brownian motion have both independent and
stationary increments.

%%%%%%%%%%%%%%%%%%%%%%%%%%%%%%%%%%%%%%%%%%%%%%%%%%%%%%%%%%%%
\subsection{Filtrations}
\label{subsec:filtrations-processes}
%%%%%%%%%%%%%%%%%%%%%%%%%%%%%%%%%%%%%%%%%%%%%%%%%%%%%%%%%%%%

A filtration is an increasing family of sigma-algebras

\[
\boxed{
\mathcal{F}_s
\subseteq
\mathcal{F}_t
\qquad
\text{whenever }s\leq t.
}
\]

The sigma-algebra

\[
\mathcal{F}_t
\]

represents information available by time \(t\).

%%%%%%%%%%%%%%%%%%%%%%%%%%%%%%%%%%%%%%%%%%%%%%%%%%%%%%%%%%%%
\subsection{Natural Filtration}
\label{subsec:natural-filtration}
%%%%%%%%%%%%%%%%%%%%%%%%%%%%%%%%%%%%%%%%%%%%%%%%%%%%%%%%%%%%

For a process \((X_t)\), its natural filtration is

\[
\boxed{
\mathcal{F}_t^X
=
\sigma(X_s:s\leq t).
}
\]

This contains exactly the information generated by the process up to
time \(t\).

%%%%%%%%%%%%%%%%%%%%%%%%%%%%%%%%%%%%%%%%%%%%%%%%%%%%%%%%%%%%
\subsection{Adapted Processes}
\label{subsec:adapted-processes}
%%%%%%%%%%%%%%%%%%%%%%%%%%%%%%%%%%%%%%%%%%%%%%%%%%%%%%%%%%%%

\begin{definitionbox}{Adapted Process}
A process \((X_t)\) is adapted to \((\mathcal{F}_t)\) if

\[
\boxed{
X_t
\text{ is }
\mathcal{F}_t\text{-measurable}
}
\]

for every \(t\).
\end{definitionbox}

Informally, the current process value is known from currently available
information.

%%%%%%%%%%%%%%%%%%%%%%%%%%%%%%%%%%%%%%%%%%%%%%%%%%%%%%%%%%%%
\subsection{Bernoulli Process}
\label{subsec:bernoulli-process}
%%%%%%%%%%%%%%%%%%%%%%%%%%%%%%%%%%%%%%%%%%%%%%%%%%%%%%%%%%%%

Let

\[
X_1,X_2,\ldots
\]

be independent Bernoulli\((p)\) random variables.

Then

\[
\boxed{
(X_n)_{n\geq1}
}
\]

is a Bernoulli process.

Each trial produces either

\[
0
\]

or

\[
1.
\]

%%%%%%%%%%%%%%%%%%%%%%%%%%%%%%%%%%%%%%%%%%%%%%%%%%%%%%%%%%%%
\subsection{Bernoulli Counting Process}
\label{subsec:bernoulli-counting-process}
%%%%%%%%%%%%%%%%%%%%%%%%%%%%%%%%%%%%%%%%%%%%%%%%%%%%%%%%%%%%

Define

\[
\boxed{
S_n
=
X_1+\cdots+X_n.
}
\]

Then \(S_n\) counts the number of successes in the first \(n\) trials.

We have

\[
\boxed{
S_n
\sim
\operatorname{Bin}(n,p).
}
\]

Also,

\[
\mathbb{E}[S_n]=np
\]

and

\[
\operatorname{Var}(S_n)=np(1-p).
\]

%%%%%%%%%%%%%%%%%%%%%%%%%%%%%%%%%%%%%%%%%%%%%%%%%%%%%%%%%%%%
\subsection{Simple Random Walk}
\label{subsec:simple-random-walk}
%%%%%%%%%%%%%%%%%%%%%%%%%%%%%%%%%%%%%%%%%%%%%%%%%%%%%%%%%%%%

Let

\[
X_1,X_2,\ldots
\]

be independent variables with

\[
P(X_i=1)=p
\]

and

\[
P(X_i=-1)=1-p.
\]

Define

\[
\boxed{
S_n
=
X_1+\cdots+X_n.
}
\]

Then

\[
(S_n)_{n\geq0}
\]

with

\[
S_0=0
\]

is a simple random walk.

%%%%%%%%%%%%%%%%%%%%%%%%%%%%%%%%%%%%%%%%%%%%%%%%%%%%%%%%%%%%
\subsection{Symmetric Random Walk}
\label{subsec:symmetric-random-walk}
%%%%%%%%%%%%%%%%%%%%%%%%%%%%%%%%%%%%%%%%%%%%%%%%%%%%%%%%%%%%

When

\[
p=\frac12,
\]

the walk is symmetric.

Then

\[
P(X_i=1)
=
P(X_i=-1)
=
\frac12.
\]

The expected increment is

\[
\mathbb{E}[X_i]=0.
\]

Therefore,

\[
\boxed{
\mathbb{E}[S_n]=0.
}
\]

%%%%%%%%%%%%%%%%%%%%%%%%%%%%%%%%%%%%%%%%%%%%%%%%%%%%%%%%%%%%
\subsection{Mean and Variance of a Random Walk}
\label{subsec:random-walk-mean-variance}
%%%%%%%%%%%%%%%%%%%%%%%%%%%%%%%%%%%%%%%%%%%%%%%%%%%%%%%%%%%%

For the general walk,

\[
\mathbb{E}[X_i]
=
p-(1-p)
=
2p-1.
\]

Therefore,

\[
\boxed{
\mathbb{E}[S_n]
=
n(2p-1).
}
\]

Since

\[
X_i^2=1,
\]

\[
\operatorname{Var}(X_i)
=
1-(2p-1)^2
=
4p(1-p).
\]

Hence,

\[
\boxed{
\operatorname{Var}(S_n)
=
4np(1-p).
}
\]

%%%%%%%%%%%%%%%%%%%%%%%%%%%%%%%%%%%%%%%%%%%%%%%%%%%%%%%%%%%%
\subsection{Worked Example: Symmetric Random Walk}
\label{subsec:worked-symmetric-random-walk}
%%%%%%%%%%%%%%%%%%%%%%%%%%%%%%%%%%%%%%%%%%%%%%%%%%%%%%%%%%%%

\begin{workedexamplebox}{Location After 100 Steps}

For a symmetric random walk,

\[
p=\frac12.
\]

Then

\[
\mathbb{E}[S_{100}]=0.
\]

Also,

\[
\operatorname{Var}(S_{100})
=
100.
\]

Therefore,

\begin{answerbox}
\[
\boxed{
\mathbb{E}[S_{100}]=0,
\qquad
\operatorname{SD}(S_{100})=10.
}
\]
\end{answerbox}

\end{workedexamplebox}

%%%%%%%%%%%%%%%%%%%%%%%%%%%%%%%%%%%%%%%%%%%%%%%%%%%%%%%%%%%%
\subsection{Distribution of a Simple Random Walk}
\label{subsec:simple-random-walk-distribution}
%%%%%%%%%%%%%%%%%%%%%%%%%%%%%%%%%%%%%%%%%%%%%%%%%%%%%%%%%%%%

Let

\[
H_n
\]

be the number of \(+1\) steps.

Then

\[
H_n\sim\operatorname{Bin}(n,p).
\]

The number of \(-1\) steps is

\[
n-H_n.
\]

Thus,

\[
S_n
=
H_n-(n-H_n)
=
2H_n-n.
\]

Therefore,

\[
\boxed{
S_n=2H_n-n.
}
\]

This expresses the random-walk position through a binomial variable.

%%%%%%%%%%%%%%%%%%%%%%%%%%%%%%%%%%%%%%%%%%%%%%%%%%%%%%%%%%%%
\subsection{Random-Walk Transition Probabilities}
\label{subsec:random-walk-transition-probabilities}
%%%%%%%%%%%%%%%%%%%%%%%%%%%%%%%%%%%%%%%%%%%%%%%%%%%%%%%%%%%%

From position \(i\),

\[
P(S_{n+1}=i+1\mid S_n=i)
=
p,
\]

while

\[
P(S_{n+1}=i-1\mid S_n=i)
=
1-p.
\]

The next-step distribution depends only on the current position.

This is an example of the Markov property.

%%%%%%%%%%%%%%%%%%%%%%%%%%%%%%%%%%%%%%%%%%%%%%%%%%%%%%%%%%%%
\subsection{Hitting Times}
\label{subsec:hitting-times}
%%%%%%%%%%%%%%%%%%%%%%%%%%%%%%%%%%%%%%%%%%%%%%%%%%%%%%%%%%%%

For a state \(a\), define the hitting time

\[
\boxed{
\tau_a
=
\inf
\{
n\geq0:S_n=a
\}.
}
\]

The Greek letter

\[
\tau
\]

is called ``tau.''

The variable \(\tau_a\) records the first time the process reaches
\(a\).

%%%%%%%%%%%%%%%%%%%%%%%%%%%%%%%%%%%%%%%%%%%%%%%%%%%%%%%%%%%%
\subsection{First-Passage Times}
\label{subsec:first-passage-times}
%%%%%%%%%%%%%%%%%%%%%%%%%%%%%%%%%%%%%%%%%%%%%%%%%%%%%%%%%%%%

More generally, for a set \(A\),

\[
\boxed{
\tau_A
=
\inf
\{
t\geq0:X_t\in A
\}.
}
\]

This is called a hitting time or first-passage time.

Such quantities are fundamental in:

\[
\text{reliability},
\]

\[
\text{finance},
\]

\[
\text{queues},
\]

and

\[
\text{diffusion theory}.
\]

%%%%%%%%%%%%%%%%%%%%%%%%%%%%%%%%%%%%%%%%%%%%%%%%%%%%%%%%%%%%
\subsection{Markov Property}
\label{subsec:markov-property}
%%%%%%%%%%%%%%%%%%%%%%%%%%%%%%%%%%%%%%%%%%%%%%%%%%%%%%%%%%%%

\begin{definitionbox}{Markov Property}
A stochastic process has the Markov property if, conditional on the
present state, the future is independent of the past.
\end{definitionbox}

For a discrete-time process,

\[
\boxed{
P
(
X_{n+1}=j
\mid
X_n=i,
X_{n-1},\ldots,X_0
)
=
P(X_{n+1}=j\mid X_n=i).
}
\]

%%%%%%%%%%%%%%%%%%%%%%%%%%%%%%%%%%%%%%%%%%%%%%%%%%%%%%%%%%%%
\subsection{Markov Chains}
\label{subsec:markov-chains}
%%%%%%%%%%%%%%%%%%%%%%%%%%%%%%%%%%%%%%%%%%%%%%%%%%%%%%%%%%%%

A discrete-time process with a discrete state space and the Markov
property is called a Markov chain.

Let the state space be

\[
S
=
\{1,2,\ldots,m\}.
\]

Define transition probabilities

\[
\boxed{
p_{ij}
=
P(X_{n+1}=j\mid X_n=i).
}
\]

For a time-homogeneous chain, these probabilities do not depend on
\(n\).

%%%%%%%%%%%%%%%%%%%%%%%%%%%%%%%%%%%%%%%%%%%%%%%%%%%%%%%%%%%%
\subsection{Transition Matrix}
\label{subsec:transition-matrix}
%%%%%%%%%%%%%%%%%%%%%%%%%%%%%%%%%%%%%%%%%%%%%%%%%%%%%%%%%%%%

The transition probabilities are arranged into the matrix

\[
\boxed{
P
=
(p_{ij}).
}
\]

For a finite state space,

\[
P
=
\begin{pmatrix}
p_{11}&p_{12}&\cdots&p_{1m}\\
p_{21}&p_{22}&\cdots&p_{2m}\\
\vdots&\vdots&\ddots&\vdots\\
p_{m1}&p_{m2}&\cdots&p_{mm}
\end{pmatrix}.
\]

Each row satisfies

\[
\boxed{
\sum_jp_{ij}=1.
}
\]

%%%%%%%%%%%%%%%%%%%%%%%%%%%%%%%%%%%%%%%%%%%%%%%%%%%%%%%%%%%%
\subsection{Worked Example: Two-State Markov Chain}
\label{subsec:worked-two-state-markov}
%%%%%%%%%%%%%%%%%%%%%%%%%%%%%%%%%%%%%%%%%%%%%%%%%%%%%%%%%%%%

\begin{workedexamplebox}{Weather Model}

Suppose the states are

\[
S=\{\text{Sunny},\text{Rainy}\}.
\]

Let

\[
P
=
\begin{pmatrix}
0.8&0.2\\
0.4&0.6
\end{pmatrix}.
\]

If today is sunny, the probability tomorrow is rainy is

\[
0.2.
\]

If today is rainy, the probability tomorrow is sunny is

\[
0.4.
\]

\begin{answerbox}
\[
\boxed{
p_{\text{Sunny,Rainy}}=0.2,
\qquad
p_{\text{Rainy,Sunny}}=0.4.
}
\]
\end{answerbox}

\end{workedexamplebox}

%%%%%%%%%%%%%%%%%%%%%%%%%%%%%%%%%%%%%%%%%%%%%%%%%%%%%%%%%%%%
\subsection{Multi-Step Transition Probabilities}
\label{subsec:multi-step-transitions}
%%%%%%%%%%%%%%%%%%%%%%%%%%%%%%%%%%%%%%%%%%%%%%%%%%%%%%%%%%%%

The \(n\)-step transition probability is

\[
\boxed{
p_{ij}^{(n)}
=
P(X_n=j\mid X_0=i).
}
\]

For a time-homogeneous finite Markov chain,

\[
\boxed{
P^{(n)}
=
P^n.
}
\]

Thus matrix multiplication computes multi-step transition probabilities.

%%%%%%%%%%%%%%%%%%%%%%%%%%%%%%%%%%%%%%%%%%%%%%%%%%%%%%%%%%%%
\subsection{Chapman--Kolmogorov Equations}
\label{subsec:chapman-kolmogorov}
%%%%%%%%%%%%%%%%%%%%%%%%%%%%%%%%%%%%%%%%%%%%%%%%%%%%%%%%%%%%

For integers \(m,n\geq0\),

\[
\boxed{
p_{ij}^{(m+n)}
=
\sum_k
p_{ik}^{(m)}
p_{kj}^{(n)}.
}
\]

In matrix form,

\[
\boxed{
P^{m+n}
=
P^mP^n.
}
\]

These are the Chapman--Kolmogorov equations.

%%%%%%%%%%%%%%%%%%%%%%%%%%%%%%%%%%%%%%%%%%%%%%%%%%%%%%%%%%%%
\subsection{State Distribution}
\label{subsec:markov-state-distribution}
%%%%%%%%%%%%%%%%%%%%%%%%%%%%%%%%%%%%%%%%%%%%%%%%%%%%%%%%%%%%

Suppose the initial distribution is the row vector

\[
\boldsymbol{\pi}_0.
\]

Then after one step,

\[
\boxed{
\boldsymbol{\pi}_1
=
\boldsymbol{\pi}_0P.
}
\]

After \(n\) steps,

\[
\boxed{
\boldsymbol{\pi}_n
=
\boldsymbol{\pi}_0P^n.
}
\]

%%%%%%%%%%%%%%%%%%%%%%%%%%%%%%%%%%%%%%%%%%%%%%%%%%%%%%%%%%%%
\subsection{Stationary Distribution}
\label{subsec:stationary-distribution-markov}
%%%%%%%%%%%%%%%%%%%%%%%%%%%%%%%%%%%%%%%%%%%%%%%%%%%%%%%%%%%%

\begin{definitionbox}{Stationary Distribution}
A probability vector \(\boldsymbol{\pi}\) is stationary for transition
matrix \(P\) if

\[
\boxed{
\boldsymbol{\pi}P
=
\boldsymbol{\pi}.
}
\]

Also,

\[
\boxed{
\sum_i\pi_i=1,
\qquad
\pi_i\geq0.
}
\]
\end{definitionbox}

If the chain starts in \(\boldsymbol{\pi}\), its state distribution
remains unchanged over time.

%%%%%%%%%%%%%%%%%%%%%%%%%%%%%%%%%%%%%%%%%%%%%%%%%%%%%%%%%%%%
\subsection{Worked Example: Stationary Distribution}
\label{subsec:worked-stationary-distribution}
%%%%%%%%%%%%%%%%%%%%%%%%%%%%%%%%%%%%%%%%%%%%%%%%%%%%%%%%%%%%

\begin{workedexamplebox}{Two-State Stationary Distribution}

Consider

\[
P
=
\begin{pmatrix}
0.8&0.2\\
0.4&0.6
\end{pmatrix}.
\]

Let

\[
\boldsymbol{\pi}
=
(\pi_1,\pi_2).
\]

Solve

\[
\boldsymbol{\pi}P
=
\boldsymbol{\pi}.
\]

The first coordinate gives

\[
0.8\pi_1+0.4\pi_2
=
\pi_1.
\]

Thus,

\[
0.4\pi_2
=
0.2\pi_1,
\]

so

\[
\pi_1=2\pi_2.
\]

Using

\[
\pi_1+\pi_2=1,
\]

we obtain

\[
\pi_2=\frac13,
\qquad
\pi_1=\frac23.
\]

\begin{answerbox}
\[
\boxed{
\boldsymbol{\pi}
=
\left(
\frac23,
\frac13
\right).
}
\]
\end{answerbox}

\end{workedexamplebox}

%%%%%%%%%%%%%%%%%%%%%%%%%%%%%%%%%%%%%%%%%%%%%%%%%%%%%%%%%%%%
\subsection{Communicating States}
\label{subsec:communicating-states}
%%%%%%%%%%%%%%%%%%%%%%%%%%%%%%%%%%%%%%%%%%%%%%%%%%%%%%%%%%%%

State \(i\) can reach state \(j\) if

\[
p_{ij}^{(n)}>0
\]

for some \(n\).

If \(i\) can reach \(j\) and \(j\) can reach \(i\), the states
communicate.

One writes

\[
\boxed{
i\leftrightarrow j.
}
\]

Communication partitions the state space into communicating classes.

%%%%%%%%%%%%%%%%%%%%%%%%%%%%%%%%%%%%%%%%%%%%%%%%%%%%%%%%%%%%
\subsection{Irreducible Markov Chains}
\label{subsec:irreducible-markov-chain}
%%%%%%%%%%%%%%%%%%%%%%%%%%%%%%%%%%%%%%%%%%%%%%%%%%%%%%%%%%%%

A Markov chain is irreducible if every pair of states communicates.

Thus there is a single communicating class.

Irreducibility means that no part of the state space is permanently
isolated from another.

%%%%%%%%%%%%%%%%%%%%%%%%%%%%%%%%%%%%%%%%%%%%%%%%%%%%%%%%%%%%
\subsection{Recurrence and Transience}
\label{subsec:recurrence-transience}
%%%%%%%%%%%%%%%%%%%%%%%%%%%%%%%%%%%%%%%%%%%%%%%%%%%%%%%%%%%%

A state is recurrent if, starting from that state, the chain returns to
it with probability \(1\).

A state is transient if there is positive probability that the process
never returns.

Thus:

\[
\boxed{
\text{recurrent}
=
\text{return eventually with probability }1,
}
\]

while

\[
\boxed{
\text{transient}
=
\text{possible permanent departure}.
}
\]

%%%%%%%%%%%%%%%%%%%%%%%%%%%%%%%%%%%%%%%%%%%%%%%%%%%%%%%%%%%%
\subsection{Periodicity}
\label{subsec:markov-periodicity}
%%%%%%%%%%%%%%%%%%%%%%%%%%%%%%%%%%%%%%%%%%%%%%%%%%%%%%%%%%%%

The period of a state \(i\) is

\[
\boxed{
d(i)
=
\gcd
\{
n\geq1:
p_{ii}^{(n)}>0
\}.
}
\]

If

\[
d(i)=1,
\]

the state is aperiodic.

For irreducible chains, all states have the same period.

%%%%%%%%%%%%%%%%%%%%%%%%%%%%%%%%%%%%%%%%%%%%%%%%%%%%%%%%%%%%
\subsection{Long-Run Markov Behavior Preview}
\label{subsec:markov-long-run-behavior}
%%%%%%%%%%%%%%%%%%%%%%%%%%%%%%%%%%%%%%%%%%%%%%%%%%%%%%%%%%%%

For a finite irreducible aperiodic Markov chain, there is a unique
stationary distribution \(\boldsymbol{\pi}\), and under standard
conditions,

\[
\boxed{
P^n
\to
\begin{pmatrix}
\boldsymbol{\pi}\\
\boldsymbol{\pi}\\
\vdots
\end{pmatrix}.
}
\]

Thus long-run state probabilities become independent of the starting
state.

%%%%%%%%%%%%%%%%%%%%%%%%%%%%%%%%%%%%%%%%%%%%%%%%%%%%%%%%%%%%
\subsection{Detailed Balance}
\label{subsec:detailed-balance}
%%%%%%%%%%%%%%%%%%%%%%%%%%%%%%%%%%%%%%%%%%%%%%%%%%%%%%%%%%%%

A probability vector \(\boldsymbol{\pi}\) satisfies detailed balance if

\[
\boxed{
\pi_ip_{ij}
=
\pi_jp_{ji}
}
\]

for all states \(i,j\).

If detailed balance holds, then \(\boldsymbol{\pi}\) is stationary.

%%%%%%%%%%%%%%%%%%%%%%%%%%%%%%%%%%%%%%%%%%%%%%%%%%%%%%%%%%%%
\subsection{Reversible Markov Chains}
\label{subsec:reversible-markov-chains}
%%%%%%%%%%%%%%%%%%%%%%%%%%%%%%%%%%%%%%%%%%%%%%%%%%%%%%%%%%%%

A Markov chain is reversible with respect to \(\boldsymbol{\pi}\) if

\[
\boxed{
\pi_ip_{ij}
=
\pi_jp_{ji}.
}
\]

Reversibility means the stationary process looks probabilistically the
same forward and backward in time.

This property is important in Markov chain Monte Carlo.

%%%%%%%%%%%%%%%%%%%%%%%%%%%%%%%%%%%%%%%%%%%%%%%%%%%%%%%%%%%%
\subsection{Counting Processes}
\label{subsec:counting-processes}
%%%%%%%%%%%%%%%%%%%%%%%%%%%%%%%%%%%%%%%%%%%%%%%%%%%%%%%%%%%%

A counting process

\[
(N(t))_{t\geq0}
\]

typically satisfies:

\[
\boxed{
N(0)=0,
}
\]

\[
\boxed{
N(t)\in\{0,1,2,\ldots\},
}
\]

and

\[
\boxed{
N(t)
\text{ is nondecreasing}.
}
\]

It counts the number of events observed by time \(t\).

%%%%%%%%%%%%%%%%%%%%%%%%%%%%%%%%%%%%%%%%%%%%%%%%%%%%%%%%%%%%
\subsection{Poisson Process}
\label{subsec:poisson-process}
%%%%%%%%%%%%%%%%%%%%%%%%%%%%%%%%%%%%%%%%%%%%%%%%%%%%%%%%%%%%

\begin{definitionbox}{Poisson Process}
A counting process \((N(t))_{t\geq0}\) is a Poisson process with rate
\(\lambda>0\) if:

\begin{itemize}
    \item \(N(0)=0\);
    \item it has independent increments;
    \item it has stationary increments;
    \item for \(0\leq s<t\),
    \[
    \boxed{
    N(t)-N(s)
    \sim
    \operatorname{Poisson}
    (
    \lambda(t-s)
    ).
    }
    \]
\end{itemize}
\end{definitionbox}

%%%%%%%%%%%%%%%%%%%%%%%%%%%%%%%%%%%%%%%%%%%%%%%%%%%%%%%%%%%%
\subsection{Poisson Process Distribution}
\label{subsec:poisson-process-distribution}
%%%%%%%%%%%%%%%%%%%%%%%%%%%%%%%%%%%%%%%%%%%%%%%%%%%%%%%%%%%%

Setting

\[
s=0,
\]

we obtain

\[
\boxed{
N(t)
\sim
\operatorname{Poisson}(\lambda t).
}
\]

Therefore,

\[
\boxed{
P(N(t)=k)
=
e^{-\lambda t}
\frac{
(\lambda t)^k
}{
k!
}.
}
\]

%%%%%%%%%%%%%%%%%%%%%%%%%%%%%%%%%%%%%%%%%%%%%%%%%%%%%%%%%%%%
\subsection{Mean and Variance of a Poisson Process}
\label{subsec:poisson-process-mean-variance}
%%%%%%%%%%%%%%%%%%%%%%%%%%%%%%%%%%%%%%%%%%%%%%%%%%%%%%%%%%%%

Because

\[
N(t)
\sim
\operatorname{Poisson}(\lambda t),
\]

we have

\[
\boxed{
\mathbb{E}[N(t)]
=
\lambda t.
}
\]

Also,

\[
\boxed{
\operatorname{Var}(N(t))
=
\lambda t.
}
\]

Thus expected count grows linearly with time.

%%%%%%%%%%%%%%%%%%%%%%%%%%%%%%%%%%%%%%%%%%%%%%%%%%%%%%%%%%%%
\subsection{Worked Example: Poisson Arrivals}
\label{subsec:worked-poisson-arrivals}
%%%%%%%%%%%%%%%%%%%%%%%%%%%%%%%%%%%%%%%%%%%%%%%%%%%%%%%%%%%%

\begin{workedexamplebox}{Events in Thirty Minutes}

Suppose events occur according to a Poisson process with rate

\[
\lambda=4
\]

per hour.

In half an hour,

\[
\lambda t
=
4\left(\frac12\right)
=
2.
\]

Thus,

\[
N(1/2)
\sim
\operatorname{Poisson}(2).
\]

The probability of exactly three events is

\[
\begin{answerbox}
\[
\boxed{
P(N(1/2)=3)
=
e^{-2}\frac{2^3}{3!}.
}
\]
\end{answerbox}

\end{workedexamplebox}

%%%%%%%%%%%%%%%%%%%%%%%%%%%%%%%%%%%%%%%%%%%%%%%%%%%%%%%%%%%%
\subsection{Poisson Increments}
\label{subsec:poisson-increments}
%%%%%%%%%%%%%%%%%%%%%%%%%%%%%%%%%%%%%%%%%%%%%%%%%%%%%%%%%%%%

For

\[
0\leq s<t,
\]

the increment

\[
N(t)-N(s)
\]

counts events occurring during

\[
(s,t].
\]

Its distribution is

\[
\boxed{
N(t)-N(s)
\sim
\operatorname{Poisson}
(
\lambda(t-s)
).
}
\]

Only the interval length matters.

%%%%%%%%%%%%%%%%%%%%%%%%%%%%%%%%%%%%%%%%%%%%%%%%%%%%%%%%%%%%
\subsection{Interarrival Times}
\label{subsec:poisson-interarrival-times}
%%%%%%%%%%%%%%%%%%%%%%%%%%%%%%%%%%%%%%%%%%%%%%%%%%%%%%%%%%%%

Let

\[
T_1
\]

be the time until the first event.

Then

\[
\{T_1>t\}
=
\{N(t)=0\}.
\]

Therefore,

\[
P(T_1>t)
=
e^{-\lambda t}.
\]

Hence,

\[
\boxed{
T_1
\sim
\operatorname{Exp}(\lambda).
}
\]

%%%%%%%%%%%%%%%%%%%%%%%%%%%%%%%%%%%%%%%%%%%%%%%%%%%%%%%%%%%%
\subsection{Successive Interarrival Times}
\label{subsec:successive-interarrival-times}
%%%%%%%%%%%%%%%%%%%%%%%%%%%%%%%%%%%%%%%%%%%%%%%%%%%%%%%%%%%%

For a Poisson process, the successive interarrival times

\[
W_1,W_2,\ldots
\]

are independent and identically distributed with

\[
\boxed{
W_i
\sim
\operatorname{Exp}(\lambda).
}
\]

Thus a Poisson process can be described equivalently through:

\[
\boxed{
\text{Poisson event counts}
}
\]

or

\[
\boxed{
\text{exponential waiting times}.
}
\]

%%%%%%%%%%%%%%%%%%%%%%%%%%%%%%%%%%%%%%%%%%%%%%%%%%%%%%%%%%%%
\subsection{Arrival Times}
\label{subsec:poisson-arrival-times}
%%%%%%%%%%%%%%%%%%%%%%%%%%%%%%%%%%%%%%%%%%%%%%%%%%%%%%%%%%%%

The time of the \(n\)-th event is

\[
\boxed{
T_n
=
W_1+\cdots+W_n.
}
\]

Since the \(W_i\) are independent exponential variables,

\[
\boxed{
T_n
\sim
\operatorname{Gamma}(n,\lambda)
}
\]

under the shape-rate convention.

For integer shape, this is also called an Erlang distribution.

%%%%%%%%%%%%%%%%%%%%%%%%%%%%%%%%%%%%%%%%%%%%%%%%%%%%%%%%%%%%
\subsection{Worked Example: Time Until Third Event}
\label{subsec:worked-third-arrival}
%%%%%%%%%%%%%%%%%%%%%%%%%%%%%%%%%%%%%%%%%%%%%%%%%%%%%%%%%%%%

\begin{workedexamplebox}{Third Arrival Time}

Suppose

\[
N(t)
\]

is a Poisson process with rate \(\lambda\).

Let \(T_3\) be the time of the third arrival.

Then

\[
T_3
=
W_1+W_2+W_3
\]

where the \(W_i\) are i.i.d.

\[
\operatorname{Exp}(\lambda).
\]

Therefore,

\begin{answerbox}
\[
\boxed{
T_3
\sim
\operatorname{Gamma}(3,\lambda).
}
\]
\end{answerbox}

\end{workedexamplebox}

%%%%%%%%%%%%%%%%%%%%%%%%%%%%%%%%%%%%%%%%%%%%%%%%%%%%%%%%%%%%
\subsection{Poisson Process Memorylessness}
\label{subsec:poisson-memorylessness}
%%%%%%%%%%%%%%%%%%%%%%%%%%%%%%%%%%%%%%%%%%%%%%%%%%%%%%%%%%%%

Because exponential waiting times are memoryless,

\[
P(W>s+t\mid W>s)
=
P(W>t).
\]

Thus if no event has occurred for \(s\) units of time, the remaining
waiting time has the same exponential distribution as before.

This is one reason the Poisson process has the Markov property.

%%%%%%%%%%%%%%%%%%%%%%%%%%%%%%%%%%%%%%%%%%%%%%%%%%%%%%%%%%%%
\subsection{Conditional Arrival Times Preview}
\label{subsec:conditional-arrival-times-poisson}
%%%%%%%%%%%%%%%%%%%%%%%%%%%%%%%%%%%%%%%%%%%%%%%%%%%%%%%%%%%%

Condition on

\[
N(t)=n.
\]

Given that exactly \(n\) events occurred by time \(t\), the unordered
arrival times behave like \(n\) independent

\[
\operatorname{Uniform}(0,t)
\]

variables.

The ordered arrival times therefore behave like the order statistics of
those uniform variables.

This creates a deep connection between Poisson processes and order
statistics.

%%%%%%%%%%%%%%%%%%%%%%%%%%%%%%%%%%%%%%%%%%%%%%%%%%%%%%%%%%%%
\subsection{Superposition of Poisson Processes}
\label{subsec:superposition-poisson}
%%%%%%%%%%%%%%%%%%%%%%%%%%%%%%%%%%%%%%%%%%%%%%%%%%%%%%%%%%%%

Suppose

\[
N_1(t)
\]

and

\[
N_2(t)
\]

are independent Poisson processes with rates

\[
\lambda_1
\]

and

\[
\lambda_2.
\]

Then

\[
\boxed{
N(t)
=
N_1(t)+N_2(t)
}
\]

is a Poisson process with rate

\[
\boxed{
\lambda_1+\lambda_2.
}
\]

%%%%%%%%%%%%%%%%%%%%%%%%%%%%%%%%%%%%%%%%%%%%%%%%%%%%%%%%%%%%
\subsection{Poisson Thinning}
\label{subsec:poisson-thinning}
%%%%%%%%%%%%%%%%%%%%%%%%%%%%%%%%%%%%%%%%%%%%%%%%%%%%%%%%%%%%

Suppose each event of a Poisson process with rate \(\lambda\) is
independently retained with probability \(p\).

Then the retained events form a Poisson process with rate

\[
\boxed{
p\lambda.
}
\]

The discarded events form an independent Poisson process with rate

\[
\boxed{
(1-p)\lambda.
}
\]

This is called thinning.

%%%%%%%%%%%%%%%%%%%%%%%%%%%%%%%%%%%%%%%%%%%%%%%%%%%%%%%%%%%%
\subsection{Renewal Processes}
\label{subsec:renewal-processes}
%%%%%%%%%%%%%%%%%%%%%%%%%%%%%%%%%%%%%%%%%%%%%%%%%%%%%%%%%%%%

A renewal process generalizes the Poisson process.

Let

\[
W_1,W_2,\ldots
\]

be i.i.d. positive waiting times.

Define arrival times

\[
\boxed{
T_n
=
W_1+\cdots+W_n.
}
\]

Then define

\[
\boxed{
N(t)
=
\max
\{
n:T_n\leq t
\}.
}
\]

The process \(N(t)\) counts renewals by time \(t\).

%%%%%%%%%%%%%%%%%%%%%%%%%%%%%%%%%%%%%%%%%%%%%%%%%%%%%%%%%%%%
\subsection{Poisson Process as a Renewal Process}
\label{subsec:poisson-renewal-special-case}
%%%%%%%%%%%%%%%%%%%%%%%%%%%%%%%%%%%%%%%%%%%%%%%%%%%%%%%%%%%%

If the waiting times satisfy

\[
W_i
\sim
\operatorname{Exp}(\lambda),
\]

then the resulting renewal process is exactly a Poisson process.

Thus,

\[
\boxed{
\text{Poisson process}
=
\text{renewal process with exponential interarrivals}.
}
\]

%%%%%%%%%%%%%%%%%%%%%%%%%%%%%%%%%%%%%%%%%%%%%%%%%%%%%%%%%%%%
\subsection{Renewal Function Preview}
\label{subsec:renewal-function-preview}
%%%%%%%%%%%%%%%%%%%%%%%%%%%%%%%%%%%%%%%%%%%%%%%%%%%%%%%%%%%%

The expected number of renewals by time \(t\) is called the renewal
function:

\[
\boxed{
m(t)
=
\mathbb{E}[N(t)].
}
\]

For a Poisson process,

\[
m(t)=\lambda t.
\]

For general renewal processes, \(m(t)\) satisfies a renewal equation.

%%%%%%%%%%%%%%%%%%%%%%%%%%%%%%%%%%%%%%%%%%%%%%%%%%%%%%%%%%%%
\subsection{Elementary Renewal Theorem Preview}
\label{subsec:elementary-renewal-theorem-preview}
%%%%%%%%%%%%%%%%%%%%%%%%%%%%%%%%%%%%%%%%%%%%%%%%%%%%%%%%%%%%

If the interarrival times have finite mean

\[
\mu_W
=
\mathbb{E}[W_1],
\]

then under suitable conditions,

\[
\boxed{
\frac{
\mathbb{E}[N(t)]
}{t}
\to
\frac1{\mu_W}.
}
\]

A stronger pathwise form gives

\[
\boxed{
\frac{N(t)}{t}
\to
\frac1{\mu_W}
}
\]

almost surely under standard assumptions.

Thus long-run renewal rate is reciprocal mean interarrival time.

%%%%%%%%%%%%%%%%%%%%%%%%%%%%%%%%%%%%%%%%%%%%%%%%%%%%%%%%%%%%
\subsection{Branching Processes}
\label{subsec:branching-processes}
%%%%%%%%%%%%%%%%%%%%%%%%%%%%%%%%%%%%%%%%%%%%%%%%%%%%%%%%%%%%

A branching process models a population in which individuals reproduce
randomly.

Let

\[
Z_n
\]

be the population size in generation \(n\).

Each individual independently produces a random number of offspring with
common distribution

\[
X.
\]

Then

\[
\boxed{
Z_{n+1}
=
\sum_{i=1}^{Z_n}
X_{n,i}.
}
\]

This is the Galton--Watson branching process.

%%%%%%%%%%%%%%%%%%%%%%%%%%%%%%%%%%%%%%%%%%%%%%%%%%%%%%%%%%%%
\subsection{Branching Process Mean}
\label{subsec:branching-process-mean}
%%%%%%%%%%%%%%%%%%%%%%%%%%%%%%%%%%%%%%%%%%%%%%%%%%%%%%%%%%%%

Let

\[
m
=
\mathbb{E}[X]
\]

be the mean number of offspring per individual.

Conditioning on \(Z_n\),

\[
\mathbb{E}[Z_{n+1}\mid Z_n]
=
mZ_n.
\]

Taking expectations,

\[
\mathbb{E}[Z_{n+1}]
=
m\mathbb{E}[Z_n].
\]

If

\[
Z_0=1,
\]

then

\[
\boxed{
\mathbb{E}[Z_n]
=
m^n.
}
\]

%%%%%%%%%%%%%%%%%%%%%%%%%%%%%%%%%%%%%%%%%%%%%%%%%%%%%%%%%%%%
\subsection{Criticality in Branching Processes}
\label{subsec:branching-criticality}
%%%%%%%%%%%%%%%%%%%%%%%%%%%%%%%%%%%%%%%%%%%%%%%%%%%%%%%%%%%%

Branching processes are classified by \(m\):

\[
\boxed{
m<1
\quad
\text{subcritical},
}
\]

\[
\boxed{
m=1
\quad
\text{critical},
}
\]

\[
\boxed{
m>1
\quad
\text{supercritical}.
}
\]

This classification strongly influences extinction behavior.

%%%%%%%%%%%%%%%%%%%%%%%%%%%%%%%%%%%%%%%%%%%%%%%%%%%%%%%%%%%%
\subsection{Extinction Probability}
\label{subsec:branching-extinction}
%%%%%%%%%%%%%%%%%%%%%%%%%%%%%%%%%%%%%%%%%%%%%%%%%%%%%%%%%%%%

Let

\[
q
\]

be the probability that the population eventually becomes extinct.

Let

\[
G(s)
=
\mathbb{E}[s^X]
\]

be the offspring PGF.

Then \(q\) satisfies the fixed-point equation

\[
\boxed{
q
=
G(q).
}
\]

The extinction probability is the smallest solution in \([0,1]\).

%%%%%%%%%%%%%%%%%%%%%%%%%%%%%%%%%%%%%%%%%%%%%%%%%%%%%%%%%%%%
\subsection{Worked Example: Simple Branching Model}
\label{subsec:worked-branching-extinction}
%%%%%%%%%%%%%%%%%%%%%%%%%%%%%%%%%%%%%%%%%%%%%%%%%%%%%%%%%%%%

\begin{workedexamplebox}{Zero or Two Offspring}

Suppose each individual has

\[
0
\]

offspring with probability \(1/2\), and

\[
2
\]

offspring with probability \(1/2\).

Then

\[
G(s)
=
\frac12
+
\frac12s^2.
\]

The extinction equation is

\[
q
=
\frac12
+
\frac12q^2.
\]

Thus,

\[
q^2-2q+1=0.
\]

Therefore,

\[
(q-1)^2=0.
\]

\begin{answerbox}
\[
\boxed{
q=1.
}
\]
\end{answerbox}

This critical process becomes extinct with probability \(1\).

\end{workedexamplebox}

%%%%%%%%%%%%%%%%%%%%%%%%%%%%%%%%%%%%%%%%%%%%%%%%%%%%%%%%%%%%
\subsection{Martingales}
\label{subsec:martingales-stochastic-processes}
%%%%%%%%%%%%%%%%%%%%%%%%%%%%%%%%%%%%%%%%%%%%%%%%%%%%%%%%%%%%

\begin{definitionbox}{Martingale}
A discrete-time process

\[
(M_n)
\]

adapted to a filtration

\[
(\mathcal{F}_n)
\]

is a martingale if

\[
\mathbb{E}[|M_n|]<\infty
\]

and

\[
\boxed{
\mathbb{E}
[
M_{n+1}
\mid
\mathcal{F}_n
]
=
M_n
}
\]

almost surely.
\end{definitionbox}

A martingale represents a fair-game structure.

%%%%%%%%%%%%%%%%%%%%%%%%%%%%%%%%%%%%%%%%%%%%%%%%%%%%%%%%%%%%
\subsection{Submartingales and Supermartingales}
\label{subsec:sub-super-martingales}
%%%%%%%%%%%%%%%%%%%%%%%%%%%%%%%%%%%%%%%%%%%%%%%%%%%%%%%%%%%%

A submartingale satisfies

\[
\boxed{
\mathbb{E}
[
M_{n+1}
\mid
\mathcal{F}_n
]
\geq
M_n.
}
\]

A supermartingale satisfies

\[
\boxed{
\mathbb{E}
[
M_{n+1}
\mid
\mathcal{F}_n
]
\leq
M_n.
}
\]

Thus:

\[
\boxed{
\text{submartingale}
=
\text{nonnegative conditional drift},
}
\]

\[
\boxed{
\text{martingale}
=
\text{zero conditional drift},
}
\]

\[
\boxed{
\text{supermartingale}
=
\text{nonpositive conditional drift}.
}
\]

%%%%%%%%%%%%%%%%%%%%%%%%%%%%%%%%%%%%%%%%%%%%%%%%%%%%%%%%%%%%
\subsection{Random Walk as a Martingale}
\label{subsec:random-walk-martingale}
%%%%%%%%%%%%%%%%%%%%%%%%%%%%%%%%%%%%%%%%%%%%%%%%%%%%%%%%%%%%

Let

\[
S_n
=
X_1+\cdots+X_n
\]

where the increments satisfy

\[
\mathbb{E}[X_i]=0
\]

and are independent of the past.

Then

\[
\mathbb{E}[S_{n+1}\mid\mathcal{F}_n]
=
S_n
+
\mathbb{E}[X_{n+1}\mid\mathcal{F}_n].
\]

Independence and zero mean give

\[
\boxed{
\mathbb{E}[S_{n+1}\mid\mathcal{F}_n]
=
S_n.
}
\]

Therefore \(S_n\) is a martingale.

%%%%%%%%%%%%%%%%%%%%%%%%%%%%%%%%%%%%%%%%%%%%%%%%%%%%%%%%%%%%
\subsection{Centered Random Walk}
\label{subsec:centered-random-walk-martingale}
%%%%%%%%%%%%%%%%%%%%%%%%%%%%%%%%%%%%%%%%%%%%%%%%%%%%%%%%%%%%

If

\[
\mathbb{E}[X_i]=\mu,
\]

then

\[
S_n-n\mu
\]

has zero conditional drift.

Therefore,

\[
\boxed{
M_n
=
S_n-n\mu
}
\]

is a martingale under the usual independence assumptions.

%%%%%%%%%%%%%%%%%%%%%%%%%%%%%%%%%%%%%%%%%%%%%%%%%%%%%%%%%%%%
\subsection{Quadratic Martingale for a Random Walk}
\label{subsec:quadratic-random-walk-martingale}
%%%%%%%%%%%%%%%%%%%%%%%%%%%%%%%%%%%%%%%%%%%%%%%%%%%%%%%%%%%%

Suppose the increments have mean \(0\) and variance \(\sigma^2\).

Then another important martingale is

\[
\boxed{
M_n
=
S_n^2-n\sigma^2.
}
\]

This compensates the predictable growth of the square by subtracting its
expected variance accumulation.

%%%%%%%%%%%%%%%%%%%%%%%%%%%%%%%%%%%%%%%%%%%%%%%%%%%%%%%%%%%%
\subsection{Worked Example: Symmetric Random Walk Martingale}
\label{subsec:worked-random-walk-martingale}
%%%%%%%%%%%%%%%%%%%%%%%%%%%%%%%%%%%%%%%%%%%%%%%%%%%%%%%%%%%%

\begin{workedexamplebox}{Fair Random Walk}

For a symmetric random walk,

\[
P(X_i=1)=P(X_i=-1)=\frac12.
\]

Then

\[
\mathbb{E}[X_i]=0.
\]

Hence

\[
S_n
=
X_1+\cdots+X_n
\]

satisfies

\[
\mathbb{E}[S_{n+1}\mid\mathcal{F}_n]
=
S_n.
\]

\begin{answerbox}
\[
\boxed{
(S_n)
\text{ is a martingale}.
}
\]
\end{answerbox}

\end{workedexamplebox}

%%%%%%%%%%%%%%%%%%%%%%%%%%%%%%%%%%%%%%%%%%%%%%%%%%%%%%%%%%%%
\subsection{Stopping Times}
\label{subsec:stopping-times}
%%%%%%%%%%%%%%%%%%%%%%%%%%%%%%%%%%%%%%%%%%%%%%%%%%%%%%%%%%%%

\begin{definitionbox}{Stopping Time}
A random time \(\tau\) is a stopping time with respect to
\((\mathcal{F}_t)\) if the event

\[
\boxed{
\{\tau\leq t\}
\in
\mathcal{F}_t
}
\]

for every \(t\).
\end{definitionbox}

This means that by time \(t\), one can determine whether stopping has
already occurred.

%%%%%%%%%%%%%%%%%%%%%%%%%%%%%%%%%%%%%%%%%%%%%%%%%%%%%%%%%%%%
\subsection{Examples of Stopping Times}
\label{subsec:stopping-time-examples}
%%%%%%%%%%%%%%%%%%%%%%%%%%%%%%%%%%%%%%%%%%%%%%%%%%%%%%%%%%%%

A first hitting time

\[
\tau_a
=
\inf
\{
n:S_n=a
\}
\]

is a stopping time with respect to the natural filtration.

A fixed deterministic time

\[
\tau=10
\]

is also a stopping time.

By contrast, a time defined using future information may fail to be a
stopping time.

%%%%%%%%%%%%%%%%%%%%%%%%%%%%%%%%%%%%%%%%%%%%%%%%%%%%%%%%%%%%
\subsection{Optional Stopping Preview}
\label{subsec:optional-stopping-preview}
%%%%%%%%%%%%%%%%%%%%%%%%%%%%%%%%%%%%%%%%%%%%%%%%%%%%%%%%%%%%

Under suitable conditions, a martingale stopped at a stopping time
retains its expectation:

\[
\boxed{
\mathbb{E}[M_\tau]
=
\mathbb{E}[M_0].
}
\]

This is the basic idea behind the optional stopping theorem.

The theorem requires conditions; it cannot be applied to arbitrary
martingales and arbitrary stopping times.

%%%%%%%%%%%%%%%%%%%%%%%%%%%%%%%%%%%%%%%%%%%%%%%%%%%%%%%%%%%%
\subsection{Gambler's Ruin Preview}
\label{subsec:gamblers-ruin-preview}
%%%%%%%%%%%%%%%%%%%%%%%%%%%%%%%%%%%%%%%%%%%%%%%%%%%%%%%%%%%%

Consider a random walk on

\[
\{0,1,\ldots,N\}
\]

that stops upon reaching either

\[
0
\]

or

\[
N.
\]

Starting from \(i\), one may ask:

\[
\boxed{
P_i(\text{hit }N\text{ before }0).
}
\]

For a symmetric walk, martingale and difference-equation methods yield

\[
\boxed{
P_i(\text{hit }N\text{ before }0)
=
\frac{i}{N}.
}
\]

%%%%%%%%%%%%%%%%%%%%%%%%%%%%%%%%%%%%%%%%%%%%%%%%%%%%%%%%%%%%
\subsection{Gaussian Processes}
\label{subsec:gaussian-processes}
%%%%%%%%%%%%%%%%%%%%%%%%%%%%%%%%%%%%%%%%%%%%%%%%%%%%%%%%%%%%

\begin{definitionbox}{Gaussian Process}
A stochastic process \((X_t)\) is Gaussian if every finite-dimensional
random vector

\[
\boxed{
(X_{t_1},\ldots,X_{t_n})
}
\]

has a multivariate normal distribution.
\end{definitionbox}

A Gaussian process is completely determined by its:

\[
\boxed{
\text{mean function}
}
\]

and

\[
\boxed{
\text{covariance function}.
}
\]

%%%%%%%%%%%%%%%%%%%%%%%%%%%%%%%%%%%%%%%%%%%%%%%%%%%%%%%%%%%%
\subsection{Gaussian Process Mean and Covariance}
\label{subsec:gaussian-process-mean-covariance}
%%%%%%%%%%%%%%%%%%%%%%%%%%%%%%%%%%%%%%%%%%%%%%%%%%%%%%%%%%%%

Define

\[
m(t)
=
\mathbb{E}[X_t]
\]

and

\[
K(s,t)
=
\operatorname{Cov}(X_s,X_t).
\]

Then the pair

\[
\boxed{
(m,K)
}
\]

determines the finite-dimensional distributions of a Gaussian process.

This is a special feature of Gaussian models.

%%%%%%%%%%%%%%%%%%%%%%%%%%%%%%%%%%%%%%%%%%%%%%%%%%%%%%%%%%%%
\subsection{Brownian Motion Preview}
\label{subsec:brownian-motion-preview}
%%%%%%%%%%%%%%%%%%%%%%%%%%%%%%%%%%%%%%%%%%%%%%%%%%%%%%%%%%%%

One of the most important Gaussian processes is Brownian motion.

A standard Brownian motion

\[
(B_t)_{t\geq0}
\]

satisfies:

\[
\boxed{
B_0=0,
}
\]

has independent increments,

has stationary Gaussian increments, and has continuous sample paths
almost surely.

For

\[
0\leq s<t,
\]

\[
\boxed{
B_t-B_s
\sim
N(0,t-s).
}
\]

%%%%%%%%%%%%%%%%%%%%%%%%%%%%%%%%%%%%%%%%%%%%%%%%%%%%%%%%%%%%
\subsection{Brownian Mean and Covariance}
\label{subsec:brownian-mean-covariance}
%%%%%%%%%%%%%%%%%%%%%%%%%%%%%%%%%%%%%%%%%%%%%%%%%%%%%%%%%%%%

For standard Brownian motion,

\[
\boxed{
\mathbb{E}[B_t]=0.
}
\]

Also,

\[
\boxed{
\operatorname{Var}(B_t)=t.
}
\]

Its covariance function is

\[
\boxed{
\operatorname{Cov}(B_s,B_t)
=
\min(s,t).
}
\]

This covariance completely determines its Gaussian finite-dimensional
laws.

%%%%%%%%%%%%%%%%%%%%%%%%%%%%%%%%%%%%%%%%%%%%%%%%%%%%%%%%%%%%
\subsection{Deriving Brownian Covariance}
\label{subsec:derive-brownian-covariance}
%%%%%%%%%%%%%%%%%%%%%%%%%%%%%%%%%%%%%%%%%%%%%%%%%%%%%%%%%%%%

Suppose

\[
s\leq t.
\]

Write

\[
B_t
=
B_s+(B_t-B_s).
\]

Then

\[
\operatorname{Cov}(B_s,B_t)
=
\operatorname{Cov}(B_s,B_s)
+
\operatorname{Cov}(B_s,B_t-B_s).
\]

Independent increments give

\[
\operatorname{Cov}(B_s,B_t-B_s)=0.
\]

Therefore,

\[
\operatorname{Cov}(B_s,B_t)
=
\operatorname{Var}(B_s)
=
s.
\]

Thus,

\[
\boxed{
\operatorname{Cov}(B_s,B_t)
=
\min(s,t).
}
\]

%%%%%%%%%%%%%%%%%%%%%%%%%%%%%%%%%%%%%%%%%%%%%%%%%%%%%%%%%%%%
\subsection{Brownian Scaling}
\label{subsec:brownian-scaling}
%%%%%%%%%%%%%%%%%%%%%%%%%%%%%%%%%%%%%%%%%%%%%%%%%%%%%%%%%%%%

Brownian motion has the scaling property

\[
\boxed{
(B_{ct})_{t\geq0}
\overset{d}{=}
(\sqrt{c}\,B_t)_{t\geq0}
}
\]

for

\[
c>0.
\]

Thus time scaling by \(c\) corresponds to space scaling by
\(\sqrt{c}\).

This is called self-similarity.

%%%%%%%%%%%%%%%%%%%%%%%%%%%%%%%%%%%%%%%%%%%%%%%%%%%%%%%%%%%%
\subsection{Brownian Motion as a Random-Walk Limit}
\label{subsec:brownian-random-walk-limit}
%%%%%%%%%%%%%%%%%%%%%%%%%%%%%%%%%%%%%%%%%%%%%%%%%%%%%%%%%%%%

Brownian motion can be viewed as a scaling limit of random walks.

A typical rescaled random walk has the form

\[
\boxed{
W_n(t)
=
\frac{
S_{\lfloor nt\rfloor}
}{
\sqrt{n}
}.
}
\]

Under suitable assumptions, these processes converge in distribution to
Brownian motion in a function-space sense.

This result is known as Donsker's invariance principle.

%%%%%%%%%%%%%%%%%%%%%%%%%%%%%%%%%%%%%%%%%%%%%%%%%%%%%%%%%%%%
\subsection{Donsker's Theorem Preview}
\label{subsec:donsker-preview}
%%%%%%%%%%%%%%%%%%%%%%%%%%%%%%%%%%%%%%%%%%%%%%%%%%%%%%%%%%%%

Donsker's theorem is a process-level extension of the central limit
theorem.

The scalar CLT gives

\[
\frac{S_n}{\sqrt{n}}
\xrightarrow{d}
N(0,\sigma^2)
\]

at a single time.

Donsker's theorem gives convergence of the entire rescaled path to a
Brownian-motion path.

Thus:

\[
\boxed{
\text{CLT}
\longrightarrow
\text{one-time distributional limit},
}
\]

\[
\boxed{
\text{Donsker}
\longrightarrow
\text{process-level limit}.
}
\]

%%%%%%%%%%%%%%%%%%%%%%%%%%%%%%%%%%%%%%%%%%%%%%%%%%%%%%%%%%%%
\subsection{Continuous-Time Markov Processes Preview}
\label{subsec:continuous-time-markov-preview}
%%%%%%%%%%%%%%%%%%%%%%%%%%%%%%%%%%%%%%%%%%%%%%%%%%%%%%%%%%%%

A continuous-time Markov process satisfies a Markov property indexed by
continuous time.

Examples include:

\[
\boxed{
\text{Poisson processes},
}
\]

\[
\boxed{
\text{continuous-time Markov chains},
}
\]

and

\[
\boxed{
\text{diffusion processes}.
}
\]

Their transition laws are often described by semigroups.

%%%%%%%%%%%%%%%%%%%%%%%%%%%%%%%%%%%%%%%%%%%%%%%%%%%%%%%%%%%%
\subsection{Transition Semigroup Preview}
\label{subsec:transition-semigroup-preview}
%%%%%%%%%%%%%%%%%%%%%%%%%%%%%%%%%%%%%%%%%%%%%%%%%%%%%%%%%%%%

For a time-homogeneous Markov process, define

\[
\boxed{
P_t f(x)
=
\mathbb{E}_x[f(X_t)].
}
\]

The family \((P_t)_{t\geq0}\) satisfies the semigroup relation

\[
\boxed{
P_{s+t}
=
P_sP_t.
}
\]

This is the continuous-time analogue of

\[
P^{m+n}=P^mP^n
\]

for discrete Markov chains.

%%%%%%%%%%%%%%%%%%%%%%%%%%%%%%%%%%%%%%%%%%%%%%%%%%%%%%%%%%%%
\subsection{Continuous-Time Markov Chains Preview}
\label{subsec:ctmc-preview}
%%%%%%%%%%%%%%%%%%%%%%%%%%%%%%%%%%%%%%%%%%%%%%%%%%%%%%%%%%%%

A continuous-time Markov chain moves among discrete states at random
times.

Its behavior is described by a generator matrix

\[
\boxed{
Q=(q_{ij}).
}
\]

For \(i\neq j\),

\[
q_{ij}\geq0,
\]

while

\[
\boxed{
q_{ii}
=
-
\sum_{j\neq i}q_{ij}.
}
\]

The matrix \(Q\) determines transition dynamics.

%%%%%%%%%%%%%%%%%%%%%%%%%%%%%%%%%%%%%%%%%%%%%%%%%%%%%%%%%%%%
\subsection{Kolmogorov Equations Preview}
\label{subsec:kolmogorov-equations-preview}
%%%%%%%%%%%%%%%%%%%%%%%%%%%%%%%%%%%%%%%%%%%%%%%%%%%%%%%%%%%%

If

\[
P(t)
\]

is the transition matrix of a continuous-time Markov chain, then under
appropriate conditions it satisfies the Kolmogorov equations:

\[
\boxed{
P'(t)
=
P(t)Q
}
\]

and

\[
\boxed{
P'(t)
=
QP(t).
}
\]

These are called the forward and backward equations.

%%%%%%%%%%%%%%%%%%%%%%%%%%%%%%%%%%%%%%%%%%%%%%%%%%%%%%%%%%%%
\subsection{Birth--Death Processes Preview}
\label{subsec:birth-death-process-preview}
%%%%%%%%%%%%%%%%%%%%%%%%%%%%%%%%%%%%%%%%%%%%%%%%%%%%%%%%%%%%

A birth--death process is a continuous-time Markov chain on

\[
\{0,1,2,\ldots\}
\]

in which transitions occur only between neighboring states.

From state \(n\),

\[
n\to n+1
\]

occurs at birth rate

\[
\lambda_n,
\]

and

\[
n\to n-1
\]

occurs at death rate

\[
\mu_n.
\]

These models appear in queues and population dynamics.

%%%%%%%%%%%%%%%%%%%%%%%%%%%%%%%%%%%%%%%%%%%%%%%%%%%%%%%%%%%%
\subsection{Queueing Processes Preview}
\label{subsec:queueing-process-preview}
%%%%%%%%%%%%%%%%%%%%%%%%%%%%%%%%%%%%%%%%%%%%%%%%%%%%%%%%%%%%

A queue may be modeled by

\[
Q(t)
=
\text{number of customers in system at time }t.
\]

If arrivals are Poisson and service times are exponential, the queue
length often forms a birth--death process.

For the classical \(M/M/1\) queue:

\[
\lambda
=
\text{arrival rate},
\]

\[
\mu
=
\text{service rate}.
\]

Under

\[
\lambda<\mu,
\]

a stationary queue-length distribution exists.

%%%%%%%%%%%%%%%%%%%%%%%%%%%%%%%%%%%%%%%%%%%%%%%%%%%%%%%%%%%%
\subsection{Reliability Processes Preview}
\label{subsec:reliability-processes-preview}
%%%%%%%%%%%%%%%%%%%%%%%%%%%%%%%%%%%%%%%%%%%%%%%%%%%%%%%%%%%%

Component states may be represented by

\[
X_t
=
\begin{cases}
1,&\text{working at time }t,\\
0,&\text{failed at time }t.
\end{cases}
\]

Systems with repair can often be modeled using Markov chains.

Systems without repair are naturally described using lifetime and
renewal processes.

%%%%%%%%%%%%%%%%%%%%%%%%%%%%%%%%%%%%%%%%%%%%%%%%%%%%%%%%%%%%
\subsection{Time Series as Stochastic Processes}
\label{subsec:time-series-stochastic-process}
%%%%%%%%%%%%%%%%%%%%%%%%%%%%%%%%%%%%%%%%%%%%%%%%%%%%%%%%%%%%

A time series is a sequence

\[
\boxed{
X_1,X_2,\ldots
}
\]

observed over time.

Statistical time-series analysis studies one realization of an
underlying stochastic process.

Important features include:

\[
\text{trend},
\]

\[
\text{seasonality},
\]

\[
\text{autocorrelation},
\]

and

\[
\text{stationarity}.
\]

%%%%%%%%%%%%%%%%%%%%%%%%%%%%%%%%%%%%%%%%%%%%%%%%%%%%%%%%%%%%
\subsection{White Noise}
\label{subsec:white-noise-process}
%%%%%%%%%%%%%%%%%%%%%%%%%%%%%%%%%%%%%%%%%%%%%%%%%%%%%%%%%%%%

A white-noise process \((\varepsilon_t)\) commonly satisfies

\[
\boxed{
\mathbb{E}[\varepsilon_t]=0,
}
\]

\[
\boxed{
\operatorname{Var}(\varepsilon_t)=\sigma^2,
}
\]

and

\[
\boxed{
\operatorname{Cov}
(
\varepsilon_s,\varepsilon_t
)
=
0
\qquad
s\neq t.
}
\]

Independent white noise is stronger than merely uncorrelated white
noise.

%%%%%%%%%%%%%%%%%%%%%%%%%%%%%%%%%%%%%%%%%%%%%%%%%%%%%%%%%%%%
\subsection{Autoregressive Process Preview}
\label{subsec:autoregressive-process-preview}
%%%%%%%%%%%%%%%%%%%%%%%%%%%%%%%%%%%%%%%%%%%%%%%%%%%%%%%%%%%%

A first-order autoregressive process satisfies

\[
\boxed{
X_t
=
\phi X_{t-1}
+
\varepsilon_t.
}
\]

Here \(\phi\) is a parameter and \(\varepsilon_t\) is white noise.

When

\[
|\phi|<1,
\]

a stationary solution exists under standard assumptions.

This is called an AR(1) process.

%%%%%%%%%%%%%%%%%%%%%%%%%%%%%%%%%%%%%%%%%%%%%%%%%%%%%%%%%%%%
\subsection{AR(1) Mean and Covariance Preview}
\label{subsec:ar1-covariance-preview}
%%%%%%%%%%%%%%%%%%%%%%%%%%%%%%%%%%%%%%%%%%%%%%%%%%%%%%%%%%%%

For a stationary zero-mean AR(1) process,

\[
X_t
=
\phi X_{t-1}
+
\varepsilon_t,
\]

with white-noise variance \(\sigma_\varepsilon^2\),

\[
\boxed{
\operatorname{Var}(X_t)
=
\frac{
\sigma_\varepsilon^2
}{
1-\phi^2
}.
}
\]

The autocorrelation is

\[
\boxed{
\rho(h)
=
\phi^{|h|}.
}
\]

Thus dependence decays geometrically with lag.

%%%%%%%%%%%%%%%%%%%%%%%%%%%%%%%%%%%%%%%%%%%%%%%%%%%%%%%%%%%%
\subsection{Markov Versus Martingale}
\label{subsec:markov-versus-martingale}
%%%%%%%%%%%%%%%%%%%%%%%%%%%%%%%%%%%%%%%%%%%%%%%%%%%%%%%%%%%%

The Markov property concerns conditional distributions:

\[
\boxed{
\text{future depends on present, not full past}.
}
\]

The martingale property concerns conditional means:

\[
\boxed{
\mathbb{E}[M_{n+1}\mid\mathcal{F}_n]
=
M_n.
}
\]

A process may be:

\[
\text{Markov but not a martingale},
\]

\[
\text{martingale but not Markov},
\]

or both.

%%%%%%%%%%%%%%%%%%%%%%%%%%%%%%%%%%%%%%%%%%%%%%%%%%%%%%%%%%%%
\subsection{Stationarity Versus Independent Increments}
\label{subsec:stationarity-versus-increments}
%%%%%%%%%%%%%%%%%%%%%%%%%%%%%%%%%%%%%%%%%%%%%%%%%%%%%%%%%%%%

A stationary process has a distribution invariant under time shifts.

A process with stationary increments only requires the distribution of

\[
X_{t+h}-X_t
\]

to depend on \(h\).

These are different properties.

Brownian motion is not stationary because

\[
\operatorname{Var}(B_t)=t
\]

depends on \(t\).

However, Brownian motion has stationary increments.

%%%%%%%%%%%%%%%%%%%%%%%%%%%%%%%%%%%%%%%%%%%%%%%%%%%%%%%%%%%%
\subsection{Poisson Process Is Not Stationary}
\label{subsec:poisson-not-stationary}
%%%%%%%%%%%%%%%%%%%%%%%%%%%%%%%%%%%%%%%%%%%%%%%%%%%%%%%%%%%%

For a Poisson process,

\[
N(t)
\sim
\operatorname{Poisson}(\lambda t).
\]

Therefore the distribution of \(N(t)\) changes with \(t\).

Hence \(N(t)\) itself is not stationary.

However,

\[
N(t+h)-N(t)
\sim
\operatorname{Poisson}(\lambda h),
\]

so the increments are stationary.

%%%%%%%%%%%%%%%%%%%%%%%%%%%%%%%%%%%%%%%%%%%%%%%%%%%%%%%%%%%%
\subsection{Process-Level Law of Large Numbers}
\label{subsec:process-lln}
%%%%%%%%%%%%%%%%%%%%%%%%%%%%%%%%%%%%%%%%%%%%%%%%%%%%%%%%%%%%

For a Poisson process,

\[
\boxed{
\frac{N(t)}{t}
\xrightarrow{\text{a.s.}}
\lambda.
}
\]

For a random walk with i.i.d. increments of finite mean \(\mu\),

\[
\boxed{
\frac{S_n}{n}
\xrightarrow{\text{a.s.}}
\mu.
}
\]

These are process-level long-run rate statements.

%%%%%%%%%%%%%%%%%%%%%%%%%%%%%%%%%%%%%%%%%%%%%%%%%%%%%%%%%%%%
\subsection{Process-Level Central Limit Behavior}
\label{subsec:process-clt}
%%%%%%%%%%%%%%%%%%%%%%%%%%%%%%%%%%%%%%%%%%%%%%%%%%%%%%%%%%%%

For a Poisson process with rate \(\lambda\),

\[
\frac{
N(t)-\lambda t
}{
\sqrt{\lambda t}
}
\]

is approximately standard normal for large \(t\).

For random walks,

\[
\frac{
S_n-n\mu
}{
\sigma\sqrt{n}
}
\xrightarrow{d}
N(0,1).
\]

Thus the LLN describes long-run drift, while the CLT describes
fluctuations around that drift.

%%%%%%%%%%%%%%%%%%%%%%%%%%%%%%%%%%%%%%%%%%%%%%%%%%%%%%%%%%%%
\subsection{Stochastic Process Modeling Workflow}
\label{subsec:stochastic-process-modeling-workflow}
%%%%%%%%%%%%%%%%%%%%%%%%%%%%%%%%%%%%%%%%%%%%%%%%%%%%%%%%%%%%

When constructing a stochastic-process model, ask:

\begin{enumerate}
    \item What is the index set?
    \item Is time discrete or continuous?
    \item What is the state space?
    \item What quantities are random?
    \item What information is available at each time?
    \item Are increments independent?
    \item Are increments stationary?
    \item Does the process have the Markov property?
    \item Is stationarity reasonable?
    \item What are the mean and covariance functions?
    \item Are hitting times important?
    \item Is a counting, renewal, branching, Markov, or Gaussian model
          appropriate?
\end{enumerate}

%%%%%%%%%%%%%%%%%%%%%%%%%%%%%%%%%%%%%%%%%%%%%%%%%%%%%%%%%%%%
\subsection{Common Errors}
\label{subsec:stochastic-process-common-errors}
%%%%%%%%%%%%%%%%%%%%%%%%%%%%%%%%%%%%%%%%%%%%%%%%%%%%%%%%%%%%

\subsubsection{Confusing a Process with a Single Random Variable}

A process is a family

\[
\{X_t:t\in T\},
\]

not a single random variable.

%%%%%%%%%%%%%%%%%%%%%%%%%%%%%%%%%%%%%%%%%%%%%%%%%%%%%%%%%%%%
\subsubsection{Confusing State Space and Index Set}

The index set tells when or where the process is observed.

The state space tells what values it can take.

%%%%%%%%%%%%%%%%%%%%%%%%%%%%%%%%%%%%%%%%%%%%%%%%%%%%%%%%%%%%
\subsubsection{Confusing Stationarity with Stationary Increments}

A process may have stationary increments without being stationary.

Brownian motion and the Poisson process are important examples.

%%%%%%%%%%%%%%%%%%%%%%%%%%%%%%%%%%%%%%%%%%%%%%%%%%%%%%%%%%%%
\subsubsection{Assuming Zero Autocorrelation Means Independence}

Uncorrelated process values need not be independent unless stronger
conditions hold.

Gaussian processes are an important special case where covariance
structure determines dependence.

%%%%%%%%%%%%%%%%%%%%%%%%%%%%%%%%%%%%%%%%%%%%%%%%%%%%%%%%%%%%
\subsubsection{Assuming Every Markov Chain Has a Unique Stationary Distribution}

Uniqueness depends on structural conditions such as irreducibility.

Convergence to the stationary distribution usually also requires
aperiodicity in finite chains.

%%%%%%%%%%%%%%%%%%%%%%%%%%%%%%%%%%%%%%%%%%%%%%%%%%%%%%%%%%%%
\subsubsection{Confusing a Stationary Distribution with a Stationary Process}

A Markov chain can possess a stationary distribution without being
started in that distribution.

If its initial law equals the stationary distribution, then the
resulting chain is stationary in distribution.

%%%%%%%%%%%%%%%%%%%%%%%%%%%%%%%%%%%%%%%%%%%%%%%%%%%%%%%%%%%%
\subsubsection{Treating the Poisson Process Rate as a Probability}

The parameter

\[
\lambda
\]

is an event rate.

It is not restricted to the interval \([0,1]\).

%%%%%%%%%%%%%%%%%%%%%%%%%%%%%%%%%%%%%%%%%%%%%%%%%%%%%%%%%%%%
\subsubsection{Confusing Poisson Counts with Exponential Waiting Times}

Poisson distributions describe event counts.

Exponential distributions describe interarrival times.

They are related but not interchangeable.

%%%%%%%%%%%%%%%%%%%%%%%%%%%%%%%%%%%%%%%%%%%%%%%%%%%%%%%%%%%%
\subsubsection{Assuming Martingale Means Independent Increments}

A martingale need not have independent increments.

The martingale property concerns conditional expectation.

%%%%%%%%%%%%%%%%%%%%%%%%%%%%%%%%%%%%%%%%%%%%%%%%%%%%%%%%%%%%
\subsubsection{Applying Optional Stopping Without Conditions}

The equality

\[
\mathbb{E}[M_\tau]
=
\mathbb{E}[M_0]
\]

requires hypotheses.

Stopping a martingale does not automatically justify expectation
preservation in every setting.

%%%%%%%%%%%%%%%%%%%%%%%%%%%%%%%%%%%%%%%%%%%%%%%%%%%%%%%%%%%%
\subsubsection{Confusing Brownian Motion with Differentiable Motion}

Brownian sample paths are continuous almost surely but are highly
irregular.

In fact, standard Brownian paths are almost surely nowhere
differentiable.

This motivates stochastic calculus.

%%%%%%%%%%%%%%%%%%%%%%%%%%%%%%%%%%%%%%%%%%%%%%%%%%%%%%%%%%%%
\subsection{Applications}
\label{subsec:stochastic-process-applications}
%%%%%%%%%%%%%%%%%%%%%%%%%%%%%%%%%%%%%%%%%%%%%%%%%%%%%%%%%%%%

\begin{applicationbox}

\textbf{Queueing Theory.}
Arrival counts, service processes, and queue lengths are naturally
modeled by Poisson, renewal, and Markov processes.

\medskip

\textbf{Reliability.}
Failure, repair, replacement, and system lifetime can be represented by
renewal and Markov models.

\medskip

\textbf{Finance.}
Asset prices, interest rates, and derivative models use martingales,
Brownian motion, and diffusion processes.

\medskip

\textbf{Biology.}
Population growth, extinction, epidemics, and genetic evolution can be
modeled with branching and Markov processes.

\medskip

\textbf{Physics.}
Random motion, diffusion, noise, and particle systems are modeled by
Brownian and other stochastic processes.

\medskip

\textbf{Operations Research.}
Inventory levels, queues, demand, and service systems evolve randomly
through time.

\medskip

\textbf{Computer Science.}
Randomized algorithms, network traffic, Markov chain Monte Carlo, and
randomized state systems use stochastic-process theory.

\medskip

\textbf{Time Series.}
Economic, environmental, engineering, and scientific measurements over
time are modeled as discrete-time stochastic processes.

\end{applicationbox}

%%%%%%%%%%%%%%%%%%%%%%%%%%%%%%%%%%%%%%%%%%%%%%%%%%%%%%%%%%%%
\subsection{Exercises}
\label{subsec:stochastic-process-exercises}
%%%%%%%%%%%%%%%%%%%%%%%%%%%%%%%%%%%%%%%%%%%%%%%%%%%%%%%%%%%%

\begin{exercise}[1]
Define a stochastic process.
\end{exercise}

\begin{exercise}[1]
What is an index set?
\end{exercise}

\begin{exercise}[1]
Define a sample path.
\end{exercise}

\begin{exercise}[1]
Define a state space.
\end{exercise}

\begin{exercise}[1]
Distinguish discrete-time and continuous-time processes.
\end{exercise}

\begin{exercise}[1]
Define independent increments.
\end{exercise}

\begin{exercise}[1]
Define stationary increments.
\end{exercise}

\begin{exercise}[1]
Define a strictly stationary process.
\end{exercise}

\begin{exercise}[1]
Define a weakly stationary process.
\end{exercise}

\begin{exercise}[1]
Define the Markov property.
\end{exercise}

\begin{exercise}[1]
Define a transition matrix.
\end{exercise}

\begin{exercise}[1]
Define a stationary distribution.
\end{exercise}

\begin{exercise}[1]
Define a Poisson process.
\end{exercise}

\begin{exercise}[1]
Define a martingale.
\end{exercise}

\begin{exercise}[1]
Define a stopping time.
\end{exercise}

\begin{exercise}[2]
For a simple random walk with step probability \(p\), compute

\[
\mathbb{E}[X_i].
\]
\end{exercise}

\begin{exercise}[2]
For the same walk, compute

\[
\operatorname{Var}(X_i).
\]
\end{exercise}

\begin{exercise}[2]
Find

\[
\mathbb{E}[S_n]
\]

and

\[
\operatorname{Var}(S_n).
\]
\end{exercise}

\begin{exercise}[2]
For a symmetric random walk, find

\[
P(S_4=0).
\]
\end{exercise}

\begin{exercise}[2]
Let

\[
P
=
\begin{pmatrix}
0.7&0.3\\
0.2&0.8
\end{pmatrix}.
\]

Verify that \(P\) is a transition matrix.
\end{exercise}

\begin{exercise}[2]
Compute

\[
P^2
\]

for the previous matrix.
\end{exercise}

\begin{exercise}[2]
Find the stationary distribution of the previous chain.
\end{exercise}

\begin{exercise}[2]
Suppose

\[
N(t)
\]

is a Poisson process with rate \(5\) per hour.

Find

\[
P(N(2)=7).
\]
\end{exercise}

\begin{exercise}[2]
For the same process, find the distribution of

\[
N(5)-N(3).
\]
\end{exercise}

\begin{exercise}[2]
For the same process, find the expected waiting time until the next
event.
\end{exercise}

\begin{exercise}[2]
Find the distribution of the time until the fourth event.
\end{exercise}

\begin{exercise}[2]
Suppose a branching process has offspring PGF

\[
G(s)=0.4+0.6s^2.
\]

Write the extinction fixed-point equation.
\end{exercise}

\begin{exercise}[3]
Prove that

\[
S_n=2H_n-n
\]

for a simple random walk, where \(H_n\) is the number of \(+1\) steps.
\end{exercise}

\begin{exercise}[3]
Derive the distribution of \(S_n\) from the binomial distribution of
\(H_n\).
\end{exercise}

\begin{exercise}[3]
Prove the Chapman--Kolmogorov equations for a discrete Markov chain.
\end{exercise}

\begin{exercise}[3]
Show that if

\[
\boldsymbol{\pi}P=\boldsymbol{\pi},
\]

then

\[
\boldsymbol{\pi}P^n=\boldsymbol{\pi}
\]

for every positive integer \(n\).
\end{exercise}

\begin{exercise}[3]
Show that detailed balance implies stationarity.
\end{exercise}

\begin{exercise}[3]
Classify the states of a finite Markov chain into communicating classes
for a transition matrix of your choice.
\end{exercise}

\begin{exercise}[3]
Derive

\[
N(t)
\sim
\operatorname{Poisson}(\lambda t)
\]

from the Poisson-process increment definition.
\end{exercise}

\begin{exercise}[3]
Show that the first interarrival time of a Poisson process is
exponential.
\end{exercise}

\begin{exercise}[3]
Show that the \(n\)-th Poisson arrival time has a gamma distribution.
\end{exercise}

\begin{exercise}[3]
Prove that the superposition of two independent Poisson processes is
Poisson.
\end{exercise}

\begin{exercise}[3]
Prove the thinning property for a Poisson process.
\end{exercise}

\begin{exercise}[3]
Derive

\[
\mathbb{E}[Z_n]=m^n
\]

for a Galton--Watson process with \(Z_0=1\).
\end{exercise}

\begin{exercise}[3]
For a branching process, explain why the extinction probability
satisfies

\[
q=G(q).
\]
\end{exercise}

\begin{exercise}[3]
Show that a centered random walk is a martingale.
\end{exercise}

\begin{exercise}[3]
Show that

\[
S_n^2-n\sigma^2
\]

is a martingale for independent mean-zero increments with variance
\(\sigma^2\).
\end{exercise}

\begin{exercise}[3]
Show that the first hitting time of a state is a stopping time.
\end{exercise}

\begin{exercise}[3]
Use martingale ideas to derive the symmetric gambler's ruin probability

\[
P_i(\text{hit }N\text{ before }0)
=
\frac{i}{N}.
\]
\end{exercise}

\begin{exercise}[3]
Derive the Brownian covariance formula

\[
\operatorname{Cov}(B_s,B_t)
=
\min(s,t).
\]
\end{exercise}

\begin{exercise}[3]
Show that Brownian increments are stationary.
\end{exercise}

\begin{exercise}[3]
Explain why Brownian motion is not itself stationary.
\end{exercise}

\begin{exercise}[3]
For a stationary AR(1) process, derive

\[
\operatorname{Var}(X_t)
=
\frac{
\sigma_\varepsilon^2
}{
1-\phi^2
}.
\]
\end{exercise}

\begin{exercise}[3]
Derive

\[
\rho(h)=\phi^{|h|}
\]

for a stationary AR(1) process.
\end{exercise}

\begin{exercise}[4]
Study recurrence and transience of the one-dimensional symmetric random
walk.
\end{exercise}

\begin{exercise}[4]
Investigate gambler's ruin for an asymmetric random walk.
\end{exercise}

\begin{exercise}[4]
Study recurrence and transience classifications for irreducible Markov
chains.
\end{exercise}

\begin{exercise}[4]
Investigate convergence to stationarity for finite irreducible
aperiodic Markov chains.
\end{exercise}

\begin{exercise}[4]
Study reversible Markov chains and detailed balance.
\end{exercise}

\begin{exercise}[4]
Investigate Markov chain Monte Carlo.
\end{exercise}

\begin{exercise}[4]
Study continuous-time Markov chains and generator matrices.
\end{exercise}

\begin{exercise}[4]
Derive the Kolmogorov forward and backward equations.
\end{exercise}

\begin{exercise}[4]
Study birth--death processes and derive stationary distributions.
\end{exercise}

\begin{exercise}[4]
Investigate the \(M/M/1\) queue.
\end{exercise}

\begin{exercise}[4]
Study renewal equations and the elementary renewal theorem.
\end{exercise}

\begin{exercise}[4]
Investigate Galton--Watson extinction criteria.
\end{exercise}

\begin{exercise}[4]
Study martingale convergence theorems.
\end{exercise}

\begin{exercise}[4]
Investigate the optional stopping theorem and its hypotheses.
\end{exercise}

\begin{exercise}[4]
Study Doob's martingale inequalities.
\end{exercise}

\begin{exercise}[4]
Investigate Gaussian processes and positive-semidefinite covariance
kernels.
\end{exercise}

\begin{exercise}[4]
Study Brownian-motion path properties.
\end{exercise}

\begin{exercise}[4]
Investigate Donsker's invariance principle.
\end{exercise}

\begin{exercise}[4]
Study stationary stochastic processes and spectral representations.
\end{exercise}

%%%%%%%%%%%%%%%%%%%%%%%%%%%%%%%%%%%%%%%%%%%%%%%%%%%%%%%%%%%%
\subsection{Conceptual Summary}
\label{subsec:stochastic-process-summary}
%%%%%%%%%%%%%%%%%%%%%%%%%%%%%%%%%%%%%%%%%%%%%%%%%%%%%%%%%%%%

A stochastic process is a family

\[
\boxed{
\{X_t:t\in T\}.
}
\]

For a fixed outcome, the function

\[
\boxed{
t\mapsto X_t(\omega)
}
\]

is a sample path.

The mean function is

\[
\boxed{
m(t)
=
\mathbb{E}[X_t].
}
\]

The covariance function is

\[
\boxed{
C(s,t)
=
\operatorname{Cov}(X_s,X_t).
}
\]

A strictly stationary process satisfies time-shift invariance of all
finite-dimensional distributions.

A weakly stationary process has constant mean and covariance depending
only on lag.

A Markov chain satisfies

\[
\boxed{
P(X_{n+1}\mid X_n,\ldots,X_0)
=
P(X_{n+1}\mid X_n).
}
\]

Its transition matrix satisfies

\[
\boxed{
P^{(n)}=P^n.
}
\]

A stationary distribution satisfies

\[
\boxed{
\boldsymbol{\pi}P
=
\boldsymbol{\pi}.
}
\]

A Poisson process satisfies

\[
\boxed{
N(t)-N(s)
\sim
\operatorname{Poisson}
(
\lambda(t-s)
).
}
\]

Its interarrival times are

\[
\boxed{
\operatorname{Exp}(\lambda).
}
\]

Its \(n\)-th arrival time is

\[
\boxed{
\operatorname{Gamma}(n,\lambda).
}
\]

A martingale satisfies

\[
\boxed{
\mathbb{E}
[
M_{n+1}\mid\mathcal{F}_n
]
=
M_n.
}
\]

A Gaussian process is one whose finite-dimensional distributions are
multivariate normal.

Standard Brownian motion satisfies

\[
\boxed{
B_t-B_s
\sim
N(0,t-s)
}
\]

for \(s<t\), and

\[
\boxed{
\operatorname{Cov}(B_s,B_t)
=
\min(s,t).
}
\]

Stochastic processes therefore provide the mathematical framework for

\[
\boxed{
\text{random evolution}
+
\text{dependence through time}
+
\text{random paths}
+
\text{long-run behavior}.
}
\]

%%%%%%%%%%%%%%%%%%%%%%%%%%%%%%%%%%%%%%%%%%%%%%%%%%%%%%%%%%%%
\subsection{Section Connections}
\label{subsec:stochastic-process-connections}
%%%%%%%%%%%%%%%%%%%%%%%%%%%%%%%%%%%%%%%%%%%%%%%%%%%%%%%%%%%%

\chapterconnection{
Stochastic Processes extends the random-variable theory of the earlier
sections from individual or finite collections of variables to entire
time-indexed families. Measure-Theoretic Probability supplies the
filtrations, conditional expectations, product spaces, and measurability
needed to define these processes rigorously. The Laws of Large Numbers
and Central Limit Theorems describe long-run behavior of random walks,
counting processes, and time averages. The next section, Stochastic
Calculus, develops integration and differentiation methods for processes
such as Brownian motion. Stochastic Differential Equations then use
those tools to model continuously evolving random systems.
}

%%%%%%%%%%%%%%%%%%%%%%%%%%%%%%%%%%%%%%%%%%%%%%%%%%%%%%%%%%%%
% SECTION SOURCE TRACKING
%%%%%%%%%%%%%%%%%%%%%%%%%%%%%%%%%%%%%%%%%%%%%%%%%%%%%%%%%%%%

%------------------------------------------------------------
% 7.12 Stochastic Processes
%------------------------------------------------------------
%
% Source basis:
%   - Standard stochastic-process theory.
%   - Standard Markov-chain theory.
%   - Standard Poisson-process and renewal theory.
%   - Standard martingale and Gaussian-process foundations.
%
% Used for:
%   - stochastic-process definitions;
%   - index sets and state spaces;
%   - discrete-time and continuous-time processes;
%   - sample paths;
%   - finite-dimensional distributions;
%   - mean functions;
%   - covariance and autocorrelation functions;
%   - strict and weak stationarity;
%   - independent increments;
%   - stationary increments;
%   - filtrations;
%   - adapted processes;
%   - Bernoulli processes;
%   - random walks;
%   - hitting and first-passage times;
%   - Markov property;
%   - transition matrices;
%   - Chapman--Kolmogorov equations;
%   - stationary distributions;
%   - communication classes;
%   - recurrence and transience;
%   - periodicity;
%   - reversibility and detailed balance;
%   - counting processes;
%   - Poisson processes;
%   - exponential interarrival times;
%   - gamma arrival times;
%   - superposition and thinning;
%   - renewal processes;
%   - branching processes;
%   - extinction probabilities;
%   - martingales;
%   - submartingales and supermartingales;
%   - stopping times;
%   - optional stopping preview;
%   - gambler's ruin preview;
%   - Gaussian processes;
%   - Brownian motion preview;
%   - Brownian covariance and scaling;
%   - Donsker theorem preview;
%   - continuous-time Markov chains preview;
%   - generator matrices;
%   - birth--death processes;
%   - queueing-process preview;
%   - time-series connections;
%   - white noise;
%   - AR(1) preview;
%   - process-level LLN and CLT behavior;
%   - applications and exercises.
%
% Notes:
%   - Stochastic Calculus is developed next in Section 7.13.
%   - Brownian motion is introduced here and developed more deeply in
%     Stochastic Calculus.
%   - Stochastic Differential Equations are developed in Section 7.14.
%   - Queueing theory is revisited in Chapter 10.
%   - Time-series methods are developed in Chapter 8.
%
%%%%%%%%%%%%%%%%%%%%%%%%%%%%%%%%%%%%%%%%%%%%%%%%%%%%%%%%%%%%